\section{Conclusion}
\label{sec:conclusion}

We introduced a co-evolutionary framework for studying cybersecurity in
autonomous multi-agent systems, in which attacker and defender agents
interact repeatedly, retain experience, and adapt their strategies. The
scientific object is not a benchmark catalogue but the empirical study
of offensive and defensive dynamics, arms races, and interaction-induced
risks that isolated evaluations cannot observe.

A multi-seed one-versus-one heuristic ladder already exhibits the
phenomenon. Static attackers are detected; adaptive attackers pace their
actions, evade detection, and succeed more often; persistent matchups
keep moving through strategy space and produce behaviours (patience,
delayed attack, deception) that no script named. A complementary
language-model study shows that hybrid episodic agents reach measurable
attack success, whereas static hybrids, persistent hybrids and a larger
coder model do not, so co-evolutionary outcomes depend on adaptation
regime and control structure rather than on model size alone.

Several hypotheses remain open. Population scaling, ESR under
communication, and coalition dynamics (H3, H4 at scale, H6) require the
E3--E8 programme of Table~\ref{tab:matrix}. High-fidelity validation on
a provisioned Kubernetes range is a natural next step. Longer term, the
same loop can support prediction of strategy shifts, defensive
co-evolution as a training regime, and agent societies in which trust and
withholding are themselves security-relevant. In all of these settings the
methodological requirement is the same: \emph{do not freeze the opponent.}
